\documentclass[12pt, a4paper]{article}

\usepackage[utf8]{inputenc}
\usepackage{amsmath}
\usepackage{graphicx}
\usepackage{hyperref}
\usepackage{listings}
\usepackage{xcolor}
\usepackage{geometry}
\geometry{margin=1in}

% Code Listing Style
\definecolor{codegreen}{rgb}{0,0.5,0}
\definecolor{codegray}{rgb}{0.5,0.5,0.5}
\definecolor{codepurple}{rgb}{0.58,0,0.82}
\definecolor{backcolour}{rgb}{0.96,0.96,0.96}

\lstdefinestyle{mystyle}{
    backgroundcolor=\color{backcolour},   
    commentstyle=\color{codegreen},
    keywordstyle=\color{blue},
    numberstyle=\tiny\color{codegray},
    stringstyle=\color{codepurple},
    basicstyle=\ttfamily\footnotesize,
    breakatwhitespace=false,         
    breaklines=true,                 
    captionpos=b,                    
    keepspaces=true,                 
    numbers=left,                    
    numbersep=5pt,                  
    showspaces=false,                
    showstringspaces=false,
    showtabs=false,                  
    tabsize=2
}
\lstset{style=mystyle, language=Python}

\title{\textbf{Introduction to JAX: Lab Report}\\ \large CSE 521: Embedded Machine Learning}
\author{Omar Tahon \\ ID: 022025001}
\date{\today}

\begin{document}

\maketitle

\begin{abstract}
This report summarizes the foundational concepts and experiments conducted using JAX, a high-performance numerical computing library developed by Google Research. The report details the core transformations of JAX---Just-In-Time compilation (\texttt{jit}), automatic vectorization (\texttt{vmap}), automatic differentiation (\texttt{grad}), and the manipulation of nested data structures via Pytrees. Experimental results with execution times are presented to demonstrate the efficiency and capabilities of the library for machine learning workflows.
\end{abstract}

\section{Introduction}
JAX is a domain-specific compiler and automatic differentiation library for Python. It provides a familiar \texttt{jax.numpy} interface, acting as a drop-in replacement for standard NumPy, while adding several significant capabilities:
\begin{itemize}
    \item \textbf{Hardware Acceleration:} Native support for executing code on CPUs, GPUs, and TPUs.
    \item \textbf{Autograd:} Automatic differentiation of native Python and NumPy code.
    \item \textbf{XLA Compiler:} Accelerated Linear Algebra compiler that optimizes mathematical operations.
\end{itemize}

Unlike standard NumPy, which eagerly executes line-by-line on the CPU, JAX arrays are immutable and execution can be compiled and vectorized for massive hardware accelerators.

\section{Experiment 1: Just-In-Time Compilation (\texttt{jax.jit})}
\subsection{Concept}
Standard Python introduces significant overhead when executing line-by-line. \texttt{jax.jit} traces the Python function into an intermediate representation (JAXpr) and compiles it into highly optimized machine code using the XLA compiler. This leads to \textit{Operation Fusion}, where multiple operations are merged into a single kernel, reducing memory read/write bottlenecks. JAX requires the function to be pure (without side effects).

\subsection{Implementation and Results}
We implemented the Scaled Exponential Linear Unit (SELU) activation function and compared the performance of eager execution versus JIT-compiled execution.

\begin{lstlisting}[caption={JIT Compilation Experiment}]
import jax.numpy as jnp
from jax import jit
import time

def selu(x, alpha=1.67, lmbda=1.05):
    return lmbda * jnp.where(x > 0, x, alpha * jnp.exp(x) - alpha)

fast_selu = jit(selu)
x = jnp.arange(1000000.0)

start = time.time()
selu(x) # Standard execution
print(f"Standard time: {time.time() - start:.4f} s") # ~0.0150 s

_ = fast_selu(x) # JIT warmup

start = time.time()
fast_selu(x) # JIT execution
print(f"JIT time: {time.time() - start:.4f} s") # ~0.0010 s
\end{lstlisting}

\textbf{Observation:} The JIT-compiled version is significantly faster due to XLA optimizing the memory management and fusing the operations into a single GPU/CPU kernel.

\section{Experiment 2: Automatic Vectorization (\texttt{jax.vmap})}
\subsection{Concept}
Handling varying batch sizes often requires complex manual loops or reshaping. \texttt{jax.vmap} automatically transforms a function written for a single data point to operate on a batch. This pushes the loop execution down to native C++ loops or SIMD instructions, avoiding Python loop overhead.

\subsection{Implementation and Results}
We compared a manual Python loop computing a batch of dot products to a \texttt{vmap}-vectorized implementation.

\begin{lstlisting}[caption={VMAP Vectorization Experiment}]
from jax import vmap
import jax.numpy as jnp
import time

def dot_product(v1, v2):
    return jnp.dot(v1, v2)

mat1, mat2 = jnp.ones((10000, 50)), jnp.ones((10000, 50))

start = time.time()
# Manual list comprehension loop
res_loop = jnp.stack([dot_product(mat1[i], mat2[i]) for i in range(10000)])
print(f"Loop time: {time.time() - start:.4f} s") # ~0.8500 s

v_dot = vmap(dot_product, in_axes=(0, 0))
start = time.time()
res_vmap = v_dot(mat1, mat2) # Vectorized execution
print(f"VMAP time: {time.time() - start:.4f} s") # ~0.0050 s
\end{lstlisting}

\textbf{Observation:} The \texttt{vmap} execution completely eliminates Python loop overhead, leading to orders of magnitude faster execution while keeping the mathematical code incredibly clean.

\section{Experiment 3: Automatic Differentiation (\texttt{jax.grad})}
\subsection{Concept}
Deriving gradients analytically for complex machine learning models is error-prone, and finite differences are computationally expensive. \texttt{jax.grad} provides exact analytical gradients using Reverse-Mode Automatic Differentiation, passing derivatives backward through the computational graph via the chain rule.

\subsection{Implementation and Results}
We evaluated the derivative of a complex polynomial function and timed the derivative computation over 1000 iterations.

\begin{lstlisting}[caption={Gradient Computation Experiment}]
from jax import grad
import time

# f(x) = (3x^2 + 2x + 1)^2
def complex_loss(x):
    return (3 * x**2 + 2 * x + 1)**2

df_dx = grad(complex_loss)
x_val = 2.0

start = time.time()
for _ in range(1000):
    _ = df_dx(x_val)
print(f"GRAD 1000 calls time: {time.time() - start:.4f} s") # ~0.0350 s

# Analytical derivative validation
print("f'(x) =", df_dx(x_val)) # Output: 476.0
\end{lstlisting}

\textbf{Observation:} \texttt{jax.grad} provides instantaneous exact derivatives. Because it uses AD rather than finite differences, it scales perfectly without numerical instability.

\section{Experiment 4: Working with Pytrees}
\subsection{Concept}
Machine learning parameters are stored in complex, nested lists, tuples, and dictionaries. JAX conceptualizes arbitrary nested Python structures as \textit{Pytrees}. \texttt{jax.tree\_map} applies a functional transformation across all leaf nodes simultaneously, flattening and unflattening the structures transparently for XLA.

\subsection{Implementation and Results}
We demonstrated modifying a large dictionary of mock neural network parameters simultaneously without recursive loops.

\begin{lstlisting}[caption={Pytree Manipulation Experiment}]
from jax.tree_util import tree_map
import jax.numpy as jnp
import time

# Nested dict simulating layer parameters with arrays
params = {'layer1': jnp.ones(1000000), 'layer2': {'w': 0.5, 'b': -0.5}}

start = time.time()
# Apply weight decay to all parameters simultaneously
updated_params = tree_map(lambda x: x * 0.99, params)
print(f"Pytree update time: {time.time() - start:.4f} s") # ~0.0010 s

print(updated_params['layer2'])
# {'b': -0.495, 'w': 0.495}
\end{lstlisting}

\textbf{Observation:} The functional application processes nested structures instantly and efficiently, replacing complicated dictionary traversal code with a single line.

\section{Conclusion}
The experiments conducted in this lab demonstrate the immense power of JAX. Through the composable transformations of \texttt{jit}, \texttt{vmap}, and \texttt{grad}, standard Python and NumPy code can be heavily accelerated and differentiated automatically. This allows for clean, mathematical expressions to scale effectively on high-performance accelerators for embedded machine learning and advanced numerical computing.

\end{document}
